\section{Deriving recall without assuming the answer}
\label{sec:recall}
The work formulas require counts such as visited nodes, retrieved postings and scored rows.
They do not, by themselves, give recall. Recall depends on \emph{which} rows survive and on the
reference ranking. This section starts with an exact definition, then adds assumptions one at
a time so their consequences can be checked.

\Needspace{12\baselineskip}
\subsection{The exact measured definition}
Let $\mathcal E(q)$ be the eligible document set after the intended filter, across the whole
database. Let $\mathcal G_K(q)$ be its exact top $K$ under the chosen reference score, and let
$\mathcal R_K(q)$ be the returned results. Then
\begin{equation}
\operatorname{Recall@}K(q)=
\frac{|\mathcal G_K(q)\cap\mathcal R_K(q)|}{K}.
\label{eq:recall-exact}
\end{equation}
Use the same tie rule in both sets. Average this quantity over queries. Repeating one query
for timing does not create new independent quality observations.

If a candidate set $\mathcal D(q)$ is reranked with the \emph{same exact reference score}, every
true top-$K$ document in that set survives into the final top $K$. Therefore
\begin{equation}
\operatorname{Recall@}K(q)=\frac{|\mathcal G_K(q)\cap\mathcal D(q)|}{K}.
\label{eq:candidate-recall}
\end{equation}
A different reranking score breaks that identity. In particular, exact binary candidates may
still miss the true float-vector top $K$ or the documents relevant to a human question.

\Needspace{12\baselineskip}
\subsection{Separate routing loss from candidate-search loss}
Let $\mathcal U(q)$ be the eligible pool in the probed clusters. Define
\[
r_{\rm route}(q)=\frac{|\mathcal G_K(q)\cap\mathcal U(q)|}{K},\qquad
r_{\rm local}(q)=\frac{|\mathcal G_K(q)\cap\mathcal D(q)|}
 {|\mathcal G_K(q)\cap\mathcal U(q)|}.
\]
When the denominator is positive and $\mathcal D\subseteq\mathcal U$,
\begin{equation}
\operatorname{Recall@}K(q)=r_{\rm route}(q)r_{\rm local}(q).
\label{eq:recall-stages}
\end{equation}
If no true top-$K$ result enters the selected pool, total recall is zero. This factorization is
exact when both factors refer to the \emph{same true global neighbors}. However, multiplying
an average routing recall by an average local recall is generally wrong:
\[
\E[r_{\rm route}r_{\rm local}]
=\E[r_{\rm route}]\E[r_{\rm local}]+\Cov(r_{\rm route},r_{\rm local}).
\]
Nor is recall against a cluster's own local top $K$ automatically the same as recall for the
particular global neighbors routed into it. That distinction matters when transferring the
flat-pool branch measurements into a clustered model.

For A, local binary search is exact. For unlimited B with valid bounds, it is exact as well.
For full-key C with complete weighted ordering, it is exact. Their global recall can still
be limited by routing. Finite node budgets and short-key quotas add a second possible loss.

\Needspace{12\baselineskip}
\subsection{A deterministic error bound for a short prefix}
Split the full score into retained and omitted coordinates:
\begin{equation}
S=T_h+Z_h,\qquad T_h=\sum_{i\in I_h}q_i b_i,\qquad
Z_h=\sum_{i\notin I_h}q_i b_i.
\end{equation}
Because $b_i\in\{-1,+1\}$,
\begin{equation}
|Z_h|\leq A_{\rm tail}=\sum_{i\notin I_h}|q_i|.
\label{eq:tail-error}
\end{equation}
Two documents can reverse their prefix order if the omitted contributions overcome their
prefix score difference. A prefix gap greater than $2A_{\rm tail}$ is sufficient to preserve
that pair's order. This is a sufficient bound, not a typical-error estimate.

A more useful query-time certificate compares the best unseen prefix score plus
$A_{\rm tail}$ with the worst fully scored retained result, as in Equation~\eqref{eq:prefix-bound}.
It adapts to the actual query. If the bound is too loose, more keys must be visited. Returning
a fixed candidate count without checking that condition remains approximate.

For float truncation, the analogous score error obeys Cauchy--Schwarz:
\[
|q_{\rm tail}^{\top}x_{\rm tail}|
\leq \|q_{\rm tail}\|_2\|x_{\rm tail}\|_2.
\]
Truncation and binary quantization are separate approximations. Matryoshka training can improve
the usefulness of prefixes, but it does not set either bound to zero or supply a universal
recall percentage.

\Needspace{12\baselineskip}
\subsection{Random documents: what the independence assumption actually gives}
The following distribution models first omit routing. When routing is added, retain the same global true-neighbor event or apply Equation~\eqref{eq:recall-stages}; replacing it with a different local top-$K$ event changes the question. Assume nonzero score variance when computing a correlation.

First condition on a particular query $q$. Assume document signs are independent and equally
likely to be $-1$ or $+1$. Then
\begin{align}
\E[S\mid q]&=0,&\Var(S\mid q)&=\sum_{i=1}^{d}q_i^2,\\
\E[T_h\mid q]&=0,&\Var(T_h\mid q)&=\sum_{i\in I_h}q_i^2,\\
\Cov(T_h,S\mid q)&=\Var(T_h\mid q).&&
\end{align}
Hence the score correlation is exactly
\begin{equation}
\rho_q=\Corr(T_h,S\mid q)=
\sqrt{\frac{\sum_{i\in I_h}q_i^2}{\sum_{i=1}^{d}q_i^2}}.
\label{eq:rho-independent}
\end{equation}
High retained variance is not the same as a small worst-case tail bound. Many small omitted coordinates can have a modest sum of squares but a large sum of absolute values. A useful approximate prefix can therefore still lack a cheap exactness certificate.

The square root is important. Retaining a fraction $v$ of score variance gives correlation
$\sqrt v$, not recall $v$ and not recall $\sqrt v$.

\begin{examplebox}{Two random-data cases can have very different prefix recall}
With equal coordinate variances, a 24-coordinate prefix of a 256-coordinate query retains
roughly $24/256$ of expected score variance. If instead the first 24 coordinates are constructed
to hold 80\% of expected score variance, their score is much more strongly correlated with the
full score. The controlled measurements in Section~\ref{sec:evidence} test both cases. The
second is a deliberately favorable distribution, not the definition of Matryoshka embeddings.
\end{examplebox}

If query coordinates are random normals $q_i\sim\mathcal N(0,\sigma_i^2)$, we can choose a
variance profile $\sigma_i^2$. The ratio of \emph{expected} prefix and full variances is
$\sum_{i\in I_h}\sigma_i^2/\sum_i\sigma_i^2$. For a specific sampled query, use its actual
Equation~\eqref{eq:rho-independent}. The expectation of a ratio need not equal the ratio of
expectations.

\Needspace{12\baselineskip}
\subsection{Real documents need a covariance model, not just a label}
A covariance matrix records how coordinates vary together. Its diagonal entries are individual variances. An off-diagonal entry describes whether two coordinates tend to move together or in opposite directions. Independent, balanced signs have unit diagonal entries and zero off-diagonal entries.

For a general binary document vector with mean $\mu$ and covariance matrix $\Sigma$, let
$q_h$ be $q$ with omitted coordinates set to zero. Then
\begin{equation}
\rho_q=\frac{q_h^\top\Sigma q}
 {\sqrt{(q_h^\top\Sigma q_h)(q^\top\Sigma q)}}.
\label{eq:rho-covariance}
\end{equation}
This includes dependencies between coordinates. Conditioning on a cluster or filter can
change both $\mu$ and $\Sigma$. A covariance measured globally is not automatically the
conditional covariance inside the query's selected clusters.

A learned prefix may preserve ranking information through these dependencies even when its
raw coordinate-energy share is small. Conversely, a prefix can have considerable energy and
still fail to identify the extreme nearest neighbors. Estimating covariance also does not
prove a normal distribution in the tails.

\Needspace{12\baselineskip}
\subsection{Equal-weight mistakes lead to a binomial model}
Let $X_i=\ind[b_i\ne t_i]$. If the $X_i$ are independent with the same disagreement probability
$p$, and all relevant weights equal $a$, then a mismatch count over $h$ coordinates satisfies
\begin{equation}
M_h=\sum_{i=1}^{h}X_i\sim\operatorname{Binomial}(h,p),\qquad
\pi_h=aM_h.
\label{eq:binomial}
\end{equation}
A Hamming radius $r$ visits
\begin{equation}
V_h(r)=\sum_{u=0}^{r}\binom{h}{u}
\label{eq:hamming-volume}
\end{equation}
keys, and its expected candidate population is
\begin{equation}
N_*\,\Pr(M_h\leq r)
=N_*\sum_{u=0}^{r}\binom{h}{u}p^u(1-p)^{h-u}.
\label{eq:binomial-hits}
\end{equation}
Here $N_*$ is the relevant reference pool size; it is $N$ without filters or cluster restriction.
Equation~\eqref{eq:binomial-hits} is a candidate-count model, not yet recall.

If $p=1/2$, every key has equal probability and the expected count reduces to $N_*V_h(r)/2^h$.
Different $p_i$ with equal weights give a Poisson-binomial count. Unequal weights give a
weighted Bernoulli sum. Neither is described by one ordinary binomial with an average $p$.

\Needspace{12\baselineskip}
\subsection{Derive recall by asking which true neighbors pass}
Write $F_X(x)=\Pr(X\leq x)$ for a cumulative distribution function (CDF): the probability that $X$ is at or below a threshold.

Let $A_h$ be the prefix penalty and $B_h$ the omitted-coordinate penalty. Full penalty is
$A_h+B_h$. Under coordinate independence, the two sums are independent. Let $u$ define the
true full-score tail, and let $v$ be the candidate prefix-penalty threshold. Then
\begin{equation}
r_{\rm population}(u,v)=
\frac{\Pr(A_h\leq v,\ A_h+B_h\leq u)}{\Pr(A_h+B_h\leq u)}
=\frac{\displaystyle\int_{a\leq v}F_{B_h}(u-a)\,dF_{A_h}(a)}
 {F_{A_h+B_h}(u)}.
\label{eq:population-recall}
\end{equation}
Read the numerator as a weighted sum: for each allowed prefix penalty, ask whether the remaining coordinates fit within the full-score threshold, multiply by the probability of that prefix penalty, and add.

This is the logic missing from a candidate-count-only calculation: a retrieved document counts
toward recall only if it also belongs to the true full-score tail.

For equal weights and an integer full radius $r_d$, the formula becomes
\begin{equation}
\frac{\displaystyle\sum_{u=0}^{\min(h,r_h,r_d)}
 \binom{h}{u}p^u(1-p)^{h-u}
 F_{\operatorname{Bin}(d-h,p)}(r_d-u)}
 {F_{\operatorname{Bin}(d,p)}(r_d)}.
\label{eq:binomial-recall}
\end{equation}
This is exact for the stated population events. Choosing those radii to approximate top $K$
and a candidate budget is another step. Finite-sample ranks and score ties need care.

\begin{examplebox}{Candidate fraction is not recall}
Use all 64 six-bit codes as documents and an all-positive equal-weight query. The best seven
codes have zero or one mistake. Keep only codes matching the first two query bits: 16 codes
pass, so the candidate fraction is $16/64=25\%$. Five of the true seven pass: the ideal code
and the four codes whose one error is in an omitted coordinate. Recall is $5/7$, about 71.4\%,
not 25\%. Allowing one error in the two-bit prefix retrieves 48 codes and includes all seven.
Both numbers follow by counting the same true-neighbor event.
\end{examplebox}

\Needspace{12\baselineskip}
\subsection{Unequal query weights: count weighted subsets}
For a fixed query and lookup coordinates $I_h$, define
\begin{equation}
V_q(v)=\#\left\{S\subseteq I_h:\sum_{i\in S}|q_i|\leq v\right\}.
\label{eq:weighted-volume}
\end{equation}
This counts every attempted key through threshold $v$, whether occupied or empty. Under
independent uniform keys, expected candidates are $N_*V_q(v)/2^h$. With real key probabilities,
\begin{equation}
\E[F(v)\mid q]=N_*\sum_{k:\pi_q(k)\leq v}\Pr(\text{key}=k\mid q,\text{cluster},\text{filter}).
\label{eq:general-hits}
\end{equation}
For a clustered lookup, apply Equation~\eqref{eq:general-hits} separately with pool size $e_j$ and the conditional key distribution in cluster $j$, then sum the results.

An average bit-disagreement rate is insufficient to determine these joint probabilities.
The empty-key count $V_q$ and occupied candidate count $F$ are separate inputs to the latency
formula. A dense document distribution near the query key can make $F$ large while $V_q$
remains small; a sparse key space can do the reverse.

For small $h$, split weights in two, enumerate each half's subset sums, sort one list, and
count pairs whose sum is at most $v$. This meet-in-the-middle calculation uses roughly
$O(2^{h/2})$ memory and $O(h2^{h/2})$ sorting/search work. It is an analysis tool for counting
keys, not the measured latency of an online hash enumerator.

For quantized nonnegative integer weights $a_i$, a dynamic program computes probability mass:
\begin{equation}
f_{i}(s)=(1-p_i)f_{i-1}(s)+p_i f_{i-1}(s-a_i).
\label{eq:dp}
\end{equation}
Replacing probabilities by counts and adding both predecessor counts computes key volumes.
This is exact for the quantized-weight model. It does not silently become an exact calculation
for unquantized query weights. If score quantization incurs a uniform score error at most
$\epsilon$, a score gap greater than $2\epsilon$ is sufficient to preserve a pair's order.
Otherwise candidates near the boundary need exact scoring or a safe expanded bound.

\Needspace{12\baselineskip}
\subsection{The conditional-normal approximation}
For many small independent contributions, a joint-normal approximation for full and prefix
scores can be convenient. Subtract each score's mean and divide by its standard deviation. This standardizes them to $(Z,Y)$ with correlation $\rho$. For a true
fraction $a=K/N_*$ and selected candidate fraction $f$, set
\[
z_K=\Phi^{-1}(1-a),\qquad y_f=\Phi^{-1}(1-f),
\]
where $\Phi$ and $\phi$ are the standard normal CDF and density. The implied recall is
\begin{equation}
\widehat r=
\frac{1}{a}\int_{z_K}^{\infty}
\phi(z)\,\Phi\!\left(\frac{\rho z-y_f}{\sqrt{1-\rho^2}}\right)dz.
\label{eq:normal-recall}
\end{equation}
We average over the \emph{entire} true tail. Substituting only $z=z_K$ calculates a boundary
conditional probability, not the integral. When $\rho=0$, recall is $f$. In the limit of
identical rankings, $\rho=1$, it is $\min(1,f/a)$, subject to ties.

This formula is not valid simply because an embedding is called Matryoshka. A short weighted
binary prefix has finite, discrete support. Its extreme-score distribution can differ greatly
from a normal tail. At a billion documents, top-$K$ events may be far deeper in the tail than
any small validation sample covers. Strong overall score correlation does not certify 99\% recall.

The controlled study below fits score correlation on one query group and checks predicted
recall on different queries. Even there, the real-data normal approximation misses by several
percentage points. Such a model can generate hypotheses; it should not declare a real
99\%-recall winner.

\Needspace{12\baselineskip}
\subsection{What can be concluded about the three methods}
\begin{itemize}
\item \textbf{A:} exact local scanning has no local candidate-search loss; global loss comes
from routing, candidate/reranker score changes, and the declared filtering task.
\item \textbf{B:} exact bound-based traversal has the same binary result as A on the same
selected pool. Finite budgets require query-level measurements or a distribution model for
visited nodes and surviving true neighbors. A node count alone is not enough.
\item \textbf{C:} full keys with complete weighted ordering are exact locally. Short keys can
be accurate when their ranking preserves true neighbors and the candidate budget is sufficient.
That possibility should be measured, not dismissed. Their enumeration and posting costs remain
separate from the recall benefit.
\end{itemize}
The exact empirical definition is always available. Distribution models explain mechanisms;
validation determines where their numerical estimates are useful.
